\chapter{Newton's Method}

We have spent many chapters learning to differentiate. We can take a function,
find its derivative, read off the slope of its tangent line at any point, and use
that tangent line as the best straight-line stand-in for the curve nearby. In the
previous chapter, on L'Hopital's rule, we saw that very principle cut through
indeterminate forms: when a ratio collapsed into $\tfrac{0}{0}$, replacing the
curves by their local linear rates revealed the true limit. In this chapter, that same tangent line yields one of the most practical algorithms in mathematics.

Here is the problem. You have an equation. Not a friendly one like $2x - 6 = 0$,
where you just rearrange and read off $x = 3$. A genuinely ugly one, like
\[
  x^5 - 3x + 1 = 0, \qquad \cos x = x, \qquad e^{-x} = x^2.
\]
For these there is \emph{no formula}. There is no quadratic-formula analog, no
algebraic trick, that will hand you the answer. And yet these equations have real
solutions -- you can practically see one of them if you sketch the curves. The
question is: how do you actually \emph{find} the number, to as many decimal places
as you please, when no formula exists?

\begin{focusbox}
Given a function $f$, find a number $r$ with $f(r) = 0$ -- a \textbf{root} of $f$ --
when there is no formula for $r$. This single problem is pervasive throughout applied mathematics. Almost every
governing equation can be framed as solving $f(x) = 0$: finding an equilibrium state,
a break-even point, an orbital intersection, or an interest rate that balances an account.
Newton's method is the standard numerical algorithm for solving such equations.
\end{focusbox}

The strategy is short, and it is built entirely out of the tangent line. You
guess. Your guess is wrong, but instead of throwing it away you \emph{improve}
it: you replace the hard curve with its easy tangent line, solve the easy linear
problem exactly, and take that as your next guess. Then you do it again. And
again. The guesses come hurtling in toward the true root -- often doubling the
number of correct digits at every single step. This chapter explains why it
works, shows in pictures exactly what it is doing, lays out the handful of ways
it can fail, and accounts for its remarkable speed.

\chapterline

\section{The idea: replace the curve with its tangent}
\label{ch19:the-idea}

Start with a picture and a guess. We want a root of some function $f$ -- a
place where its graph crosses the horizontal axis. Suppose we have a current guess,
call it $x_0$. Almost certainly $x_0$ is not the root: the point $(x_0, f(x_0))$
sits above or below the axis, not on it.

Now here is the move. The curve near $x_0$ is hard to work with, but the
\emph{tangent line} at $x_0$ is trivial -- it is a straight line, and finding where
a straight line crosses the axis is grade-school algebra. So we make a bargain: we
pretend, just for one step, that the function \emph{is} its tangent line. We slide
down (or up) that tangent line until it hits the axis, and wherever it lands we
call our new, improved guess $x_1$.

\begin{intuitionbox}
The whole method in one sentence: \emph{you cannot easily find where the curve hits
the axis, so instead you find where its tangent line hits the axis, and you bet
that this is closer.} The tangent line hugs the curve near $x_0$, so if $x_0$ is
already in the neighborhood of the root, the place the tangent crosses is a much
better estimate than $x_0$ was. Repeat the bargain from the new point and the
estimates close in.
\end{intuitionbox}

Figure~\ref{ch19:fig-onestep} shows a single step. The guess $x_0$ is off to the
right of the true root $r$. We climb up to the point on the curve, draw the tangent
there, and follow it down to the axis. It lands at $x_1$, noticeably closer to $r$
than $x_0$ was.

\begin{figure}[H]
\centering
\begin{tikzpicture}
\begin{axis}[
  width=11cm, height=7cm,
  axis lines=middle,
  xlabel={$x$}, ylabel={$y$},
  xmin=-0.4, xmax=3.2, ymin=-2.2, ymax=5.5,
  xtick={\empty}, ytick={\empty},
  clip=false,
]
% f(x) = x^2 - 2, root at sqrt(2) ~ 1.4142
\addplot[curveblue, very thick, samples=100, domain=-0.3:2.6] {x^2 - 2};
\node[curveblue, anchor=south east, font=\small] at (axis cs:2.3,3.5) {$y=f(x)$};
% guess x0 = 2.6, f(x0) = 4.76, f'(x0)=5.2, tangent: y = 4.76 + 5.2(x-2.6)
% crosses axis at x1 = 2.6 - 4.76/5.2 = 1.6846
\addplot[tangentred, thick, samples=2, domain=1.0:2.95] {4.76 + 5.2*(x-2.6)};
% points
\addplot[only marks, mark=*, mark size=1.6pt, black] coordinates {(2.6,0) (2.6,4.76) (1.6846,0) (1.41421,0)};
\node[below] at (axis cs:2.6,-0.08) {$x_0$};
\node[below] at (axis cs:1.6846,-0.08) {$x_1$};
\node[above left] at (axis cs:1.41421,0.0) {$r$};
\node[above left] at (axis cs:2.55,4.76) {$(x_0,f(x_0))$};
% dashed vertical
\addplot[gray, dashed, samples=2, domain=0:1] coordinates {(2.6,0) (2.6,4.76)};
\end{axis}
\end{tikzpicture}
\caption{One step of Newton's method on $f(x)=x^2-2$, whose root is $r=\sqrt2$.
From the guess $x_0$ we ride the tangent line (red) down to the axis, landing at
$x_1$, much nearer the root $r$ than $x_0$ was.}
\label{ch19:fig-onestep}
\end{figure}

That is the entire idea. Everything else is bookkeeping: turning ``follow the
tangent to the axis'' into a formula, and then turning the formula into a loop.

\section{Deriving the iteration formula}
\label{ch19:formula}

Do the algebra once, carefully, because the formula is worth understanding rather
than memorizing.

We are standing at $x_0$. The point on the curve is $(x_0, f(x_0))$, and the slope
of the tangent there is $f'(x_0)$. The tangent line is therefore
\[
  y = f(x_0) + f'(x_0)\,(x - x_0).
\]
This is nothing but the linear approximation of $f$ at $x_0$ -- the very same
formula we built in the chapter on differentials. We want the place where this line
crosses the horizontal axis, i.e.\ where $y = 0$. Set the right-hand side to zero
and solve for $x$:
\[
  0 = f(x_0) + f'(x_0)\,(x - x_0)
  \quad\Longrightarrow\quad
  x - x_0 = -\,\frac{f(x_0)}{f'(x_0)}
  \quad\Longrightarrow\quad
  x = x_0 - \frac{f(x_0)}{f'(x_0)}.
\]
That landing point is our new guess $x_1$. Naming the general step, we get the
celebrated recurrence.

\begin{focusbox}
\textbf{Newton's iteration.} Starting from a guess $x_0$, generate a sequence by
\[
  \boxed{\,x_{n+1} = x_n - \dfrac{f(x_n)}{f'(x_n)}\,}
  \qquad n = 0, 1, 2, \ldots
\]
Each new estimate is the old one corrected by the amount $f(x_n)/f'(x_n)$. When the
$x_n$ stop changing -- when $f(x_n)$ has shrunk to essentially zero -- you have your
root.
\end{focusbox}

Look at what the correction term $f(x_n)/f'(x_n)$ is telling you. The numerator
$f(x_n)$ is how far off the axis you are -- your ``vertical error.'' The denominator
$f'(x_n)$ is the slope, the rate at which the function climbs. Dividing height by
slope gives you a \emph{horizontal} distance: ``if I am this high up and the ground
rises at this rate, I must back up this far to reach ground level.'' It is exactly
the run you get from a given rise and a given slope. The method is just rise over
slope equals run, applied over and over.

\textbf{Pause and think:} What happens if your starting guess $x_0$ happens to be
the exact root $r$? Then $f(x_0) = 0$, the correction term $f(x_0)/f'(x_0)$
vanishes, and $x_1 = x_0$. The iteration immediately stands still; an exact root
is a resting point of the method.

\begin{happeningbox}
What is really happening is that Newton's method \emph{inverts the linear
approximation}. Linear approximation answered the forward question: ``given a small
step $\Delta x$, how much does $f$ change?'' -- the answer being $\Delta f \approx
f'(x)\,\Delta x$. Newton's method runs that backwards. It asks: ``I want $f$ to
change by a specific amount, namely by $-f(x_n)$ so as to land on zero -- what step
$\Delta x$ achieves that?'' Solving $f'(x_n)\,\Delta x = -f(x_n)$ gives $\Delta x =
-f(x_n)/f'(x_n)$, and stepping by that amount is the update. So the iteration is
linear approximation solved for the input instead of the output. The reason it is
only \emph{approximate} -- the reason you must iterate rather than land on the root
in one shot -- is that the curve is not actually a straight line, so the real $f$ is
not yet zero at $x_{n+1}$. But it is closer, and from the closer point the linear
approximation is even better, so the next correction is smaller still.
\end{happeningbox}

\begin{whybox}
Why does following the tangent get you closer to the root at all? Because near a
point the tangent line and the curve agree to \emph{first order} -- they have the
same value and the same slope -- so over a short horizontal span they part company
only slowly, by an amount governed by the curvature (the second derivative). If
your guess is already reasonably near the root, the stretch between guess and root
is short, the tangent and curve are nearly indistinguishable across it, and the
tangent's crossing is therefore nearly the curve's crossing. The error you make by
using the line instead of the curve is small, and -- as we will prove
shortly -- it is small \emph{compared to the error you started with}. Each step's
leftover error is roughly the \emph{square} of the previous error. Square a small
number and you get a tiny number. That squaring is the engine of the method's
speed.
\end{whybox}

\section{A first worked example: square roots by hand}
\label{ch19:sqrt-example}

Nothing makes the method concrete like computing a number you already know. Take
$\sqrt2$. It is the positive root of
\[
  f(x) = x^2 - 2, \qquad f'(x) = 2x.
\]
Plug into the iteration:
\[
  x_{n+1} = x_n - \frac{x_n^2 - 2}{2 x_n}
          = x_n - \frac{x_n}{2} + \frac{1}{x_n}
          = \frac{x_n}{2} + \frac{1}{x_n}
          = \frac{1}{2}\!\left(x_n + \frac{2}{x_n}\right).
\]
That last form is famous in its own right -- it is the \emph{Babylonian method} for
square roots, known thousands of years before Newton: ``Average your guess with
$2$ divided by your guess.'' The ancients arrived at the exact Newton iteration
long before the underlying calculus was formulated.

Start, deliberately, with a poor guess: $x_0 = 2$.

\begin{labbox}
Carry these out yourself with a calculator; watch the digits lock in.
\begin{align*}
  x_0 &= 2 \\
  x_1 &= \tfrac12\!\left(2 + \tfrac{2}{2}\right) = \tfrac12(3) = 1.5 \\
  x_2 &= \tfrac12\!\left(1.5 + \tfrac{2}{1.5}\right) = \tfrac12(1.5 + 1.3333\ldots) = 1.41666\ldots \\
  x_3 &= \tfrac12\!\left(1.416667 + \tfrac{2}{1.416667}\right) = 1.4142156\ldots \\
  x_4 &= 1.41421356237\ldots
\end{align*}
The true value is $\sqrt2 = 1.41421356237\ldots$
\end{labbox}

Stare at the errors -- the distance from each guess to the true $\sqrt2$:
\begin{center}
\begin{tabular}{cl}
\toprule
$n$ & error $\abs{x_n - \sqrt2}$ \\
\midrule
$0$ & $0.586$ \\
$1$ & $0.0858$ \\
$2$ & $0.00245$ \\
$3$ & $0.0000021$ \\
$4$ & $0.0000000000016$ \\
\bottomrule
\end{tabular}
\end{center}

Read that column from the top down and feel what is happening. The error goes from
about half, to about a tenth, to a few thousandths, to a couple millionths, to a
\emph{trillionth}. Roughly speaking, the number of correct decimal digits is
\emph{doubling} at every step: $0$ correct, then $1$, then $2$--$3$, then $6$, then
$12$. After four steps from a lazy guess of $2$ we have $\sqrt2$ to twelve places.
That doubling of correct digits is the hallmark of \textbf{quadratic convergence},
and we will explain it precisely soon.

\begin{applicationbox}
This is not a toy. The square-root key on your calculator, and the
\texttt{sqrt} function buried in every programming language's math library, does
not consult a table -- it runs an iteration of essentially this kind, typically
Newton's method or a close relative, started from a clever first guess and run just
enough times to fill the available precision. Division itself, on some processors,
is computed by Newton's method applied to $f(x) = 1/x - a$, which finds $1/a$ using
only multiplication and subtraction -- no division at all. The arithmetic happening
inside your machine is shot through with this one idea.
\end{applicationbox}

\section{Watching convergence in pictures}
\label{ch19:pictures}

The formula is clean, but the method is at heart geometric, so watch several
steps unfold on a graph. Take a genuinely non-algebraic equation:
\[
  f(x) = x^3 - x - 2 = 0.
\]
There is no nice closed form you would want to use, but a sketch shows a single
real root a bit past $1.5$. Start at $x_0 = 2$.

\begin{figure}[H]
\centering
\begin{tikzpicture}
\begin{axis}[
  width=12cm, height=7.5cm,
  axis lines=middle,
  xlabel={$x$}, ylabel={$y$},
  xmin=1.2, xmax=2.3, ymin=-1.5, ymax=5.0,
  xtick={\empty},
  ytick={\empty},
  clip=false,
]
\addplot[curveblue, very thick, samples=100, domain=1.2:2.25] {x^3 - x - 2};
\addplot[tangentred, thick, samples=2, domain=1.5:2.05] {4 + 11*(x-2)};
\addplot[targetgreen, thick, samples=2, domain=1.45:1.75] {0.74530 + 7.03306*(x-1.63636)};
\addplot[only marks, mark=*, mark size=1.5pt, black] coordinates {(2,0)(2,4)(1.63636,0)(1.63636,0.74530)(1.53039,0)(1.52138,0)};
\addplot[gray,dashed,samples=2] coordinates {(2,0)(2,4)};
\addplot[gray,dashed,samples=2] coordinates {(1.63636,0)(1.63636,0.74530)};
\node[below] at (axis cs:2,-0.12) {$x_0$};
\node[below] at (axis cs:1.63636,-0.12) {$x_1$};
\node[below] at (axis cs:1.53039,-0.12) {$x_2$};
\node[above left] at (axis cs:1.515,0.05) {$r$};
\node[above left] at (axis cs:2,4) {$(x_0,f(x_0))$};
\end{axis}
\end{tikzpicture}
\caption{Three guesses converging on the root of $f(x)=x^3-x-2$. From $x_0=2$ the
red tangent lands at $x_1$; from $x_1$ the green tangent lands at $x_2$, already
nearly on top of the root. Each tangent is taken at a point closer to the
crossing -- the guesses pile up against the root.}
\label{ch19:fig-cubic}
\end{figure}

The numbers behind the picture:
\[
  x_0 = 2, \quad x_1 = 1.63636\ldots, \quad x_2 = 1.53039\ldots,
  \quad x_3 = 1.52144\ldots, \quad x_4 = 1.52138\ldots
\]
By $x_4$ the digits have stopped moving at the precision shown. The picture and the
table tell the same story: each step the new point clings tighter to the root,
because each tangent is taken at a point ever closer to the curve's actual
crossing, where the line and the curve are ever harder to tell apart.

\begin{walkbox}
Take this on a walk. In Figure~\ref{ch19:fig-cubic} the curve is concave up (it
bends upward) and the root is approached \emph{from the right}, each guess landing
to the left of the last but never overshooting past the root. Picture instead a
curve that is concave \emph{down} near its root, approached from the right. Sketch
a tangent. Does the iterate stay on the same side of the root, or jump across it?
Convince yourself that the curvature and which side you start on together decide
whether Newton's method marches in steadily from one side or zig-zags across the
root closing in.
\end{walkbox}

\section{When it works -- and when it fails}
\label{ch19:failures}

Newton's method is fast, but it is not foolproof, and a teacher who shows you only
the triumphs is doing you a disservice. The method comes with a fine-print
contract, and there are four classic ways to violate it. Knowing them is what
separates someone who can \emph{use} the method from someone who merely knows the
formula.

\subsection{The happy case: a good start near a simple root}

When everything goes right, two conditions hold. First, the starting guess $x_0$ is
\emph{close enough} to a root $r$. Second, the root is \emph{simple} -- meaning
$f'(r) \ne 0$, so the curve crosses the axis at a genuine slant rather than grazing
it tangentially. Under those conditions convergence is guaranteed to be quadratic:
the count of correct digits doubles at every step. That is the regime we have been
admiring.

The trouble is the phrase ``close enough.'' Newton's method is a \emph{local}
method. It promises wonderful things once you are in the neighborhood of the root,
and says almost nothing about how to get into that neighborhood in the first place.
Starting far away, anything can happen.

\subsection{Failure 1: the derivative vanishes -- a flat spot}

Look again at the formula:
\[
  x_{n+1} = x_n - \frac{f(x_n)}{f'(x_n)}.
\]
If ever $f'(x_n) = 0$, you are dividing by zero and the method blows up. Geometrically
this is obvious: where the derivative is zero the tangent line is \emph{horizontal},
and a horizontal line either never meets the axis or lies along it. There is no
``where the tangent crosses'' to step to. If your guess lands on or very near a flat
spot of $f$, the tangent is nearly horizontal and it slings your next guess off to
some wildly distant place -- possibly farther from the root than you began.

\begin{figure}[H]
\centering
\begin{tikzpicture}
\begin{axis}[
  width=11cm, height=6.2cm,
  axis lines=middle,
  xlabel={$x$}, ylabel={$y$},
  xmin=-2.8, xmax=3.0, ymin=-1.2, ymax=2.5,
  xtick={\empty}, ytick={\empty}, clip=false,
]
\addplot[curveblue, very thick, samples=100, domain=-1.4:2.6] {0.1*(x-1)^3 + 0.25*(x-1) + 0.75};
\addplot[tangentred, thick, samples=2, domain=-2.3:1.6] {0.75 + 0.25*(x-1)};
\addplot[only marks, mark=*, mark size=1.5pt, black] coordinates {(1,0.75)(1,0)(-2,0)};
\addplot[gray, dashed] coordinates {(1,0) (1,0.75)};
\node[below] at (axis cs:1,-0.08) {$x_0$};
\node[below] at (axis cs:-2,-0.08) {$x_1$};
\node[above] at (axis cs:1,0.8) {$(x_0,f(x_0))$};
\node[anchor=south east, font=\small, tangentred] at (axis cs:-0.1,0.48) {shallow tangent};
\end{axis}
\end{tikzpicture}
\caption{A near-horizontal tangent. The shallow slope at $x_0=1$ ($f'(x_0)=0.25$) forces
the tangent line to travel a large distance before hitting the axis, slinging the next
guess all the way to $x_1=-2$. If the derivative were exactly zero, the tangent would be
horizontal and the next step would be completely undefined.}
\label{ch19:fig-flat}
\end{figure}

\subsection{Failure 2: a bad start -- running off to infinity}

Even with a perfectly nice function, a careless starting point can send the
iterates marching off to infinity instead of inward to the root. The classic
demonstration is
\[
  f(x) = \arctan x,
\]
whose only root is $x = 0$. Its graph flattens out toward horizontal asymptotes
$\pm\pi/2$ as $\abs{x}$ grows, so far from the origin the slope is tiny. Start with a
guess that is too large -- say $x_0 = 2$ -- and the tangent, being nearly flat,
overshoots the root and lands \emph{farther out} on the other side. The next step
overshoots even more. The iterates spiral outward, $\abs{x_n} \to \infty$, and the
method diverges. There is a critical starting radius: inside it Newton converges to
$0$, outside it the method runs away. This is the price of being a local method --
``close enough'' was a real condition, not a formality (Figure~\ref{ch19:fig-overshoot}).

\begin{figure}[H]
\centering
\begin{tikzpicture}
\begin{axis}[
  width=11.5cm, height=6.2cm,
  axis lines=middle,
  xlabel={$x$}, ylabel={$y$},
  xmin=-3.6, xmax=4.3, ymin=-1.8, ymax=1.8,
  xtick={0}, xticklabels={$0$},
  ytick={-1.5708, 1.5708}, yticklabels={$-\frac{\pi}{2}$, $\frac{\pi}{2}$},
  clip=false,
  title style={font=\small},
]
% arctan curve
\addplot[curveblue, very thick, samples=120, domain=-3.5:3.5] {rad(atan(x))};
% Asymptotes
\addplot[gray, dashed, domain=-3.6:3.8] {pi/2};
\addplot[gray, dashed, domain=-3.6:3.8] {-pi/2};
\node[gray, font=\scriptsize, anchor=south east] at (axis cs:3.6,1.57) {$y=\frac{\pi}{2}$};
\node[gray, font=\scriptsize, anchor=north east] at (axis cs:3.6,-1.57) {$y=-\frac{\pi}{2}$};

% x0 = 1.6, f(1.6) = 1.0122, f'(1.6) = 0.2809 -> x1 = -2.003
\addplot[tangentred, thick, samples=2, domain=-2.3:2.2] {1.0122 + 0.2809*(x-1.6)};
\addplot[only marks, mark=*, mark size=1.8pt, black] coordinates {(1.6,0) (1.6,1.0122) (-2.003,0)};
\addplot[gray, dashed] coordinates {(1.6,0) (1.6,1.0122)};

% x1 = -2.003, f(-2.003) = -1.1078, f'(-2.003) = 0.1995 -> x2 = 3.55
\addplot[targetgreen, thick, dashed, samples=2, domain=-2.3:3.7] {-1.1078 + 0.1995*(x+2.003)};
\addplot[only marks, mark=*, mark size=1.8pt, black] coordinates {(-2.003,-1.1078) (3.55,0)};
\addplot[gray, dashed] coordinates {(-2.003,0) (-2.003,-1.1078)};

\node[below, font=\small] at (axis cs:1.6,-0.08) {$x_0$};
\node[above, font=\small, tangentred] at (axis cs:-2.003,0.08) {$x_1$};
\node[below, font=\small, targetgreen] at (axis cs:3.55,-0.08) {$x_2$};
\node[above right, font=\footnotesize, curveblue] at (axis cs:2.1,1.15) {$y=\arctan x$};
\end{axis}
\end{tikzpicture}
\caption{Failure through outward overshoot on $f(x)=\arctan x$.
Starting from $x_0=1.6$, the shallow tangent (red) overshoots the root at $x=0$, landing
farther away at $x_1 \approx -2.0$. The next tangent at $x_1$ (dashed green) overshoots
even farther to $x_2 \approx +3.55$. The iterates diverge to $\pm\infty$ in an expanding spiral.}
\label{ch19:fig-overshoot}
\end{figure}

\begin{trapbox}
A common and dangerous habit is to run a fixed number of Newton iterations and
\emph{trust the result without checking}. If the start was bad, your final $x_n$ may
have nothing to do with any root -- it might be drifting toward infinity, or
oscillating, or have parked itself near a different root than you wanted. Always
verify by plugging back in: confirm that $f(x_n)$ is genuinely small. The
iteration's internal confidence is no substitute for evaluating the function.
\end{trapbox}

\subsection{Failure 3: a cycle that never lands}

The most surprising failure is that the iterates can fall into a \emph{loop},
bouncing between two values forever, never approaching any root even though one
exists. The textbook example is
\[
  f(x) = x^3 - 2x + 2.
\]
Take $x_0 = 0$. Then $f(0) = 2$ and $f'(0) = -2$, so
\[
  x_1 = 0 - \frac{2}{-2} = 1.
\]
Now from $x_1 = 1$: $f(1) = 1 - 2 + 2 = 1$ and $f'(1) = 3 - 2 = 1$, so
\[
  x_2 = 1 - \frac{1}{1} = 0.
\]
We are back to $0$. The next step sends us to $1$ again, then back to $0$, forever:
$0 \to 1 \to 0 \to 1 \to \cdots$. The method has entered a two-cycle and will never
escape, even though $f$ has a perfectly good real root near $-1.77$. Starting from
the wrong place locked us into a trap with no exit.

\begin{figure}[H]
\centering
\begin{tikzpicture}
\begin{axis}[
  width=11cm, height=6.5cm,
  axis lines=middle,
  xlabel={$x$}, ylabel={$y$},
  xmin=-0.6, xmax=1.6, ymin=-0.6, ymax=3.0,
  xtick={0.5}, ytick={\empty}, clip=false,
]
\addplot[curveblue, very thick, samples=100, domain=-0.55:1.55] {x^3 - 2*x + 2};
% tangent at x0=0: y = 2 - 2x -> crosses axis at x=1
\addplot[tangentred, thick, samples=2, domain=0:1.1] {2 - 2*x};
% tangent at x1=1: f=1, slope 1: y = 1 + (x-1) = x -> crosses axis at x=0
\addplot[targetgreen, thick, samples=2, domain=-0.1:1.1] {x};
\addplot[only marks, mark=*, mark size=1.5pt, black] coordinates {(0,2)(0,0)(1,1)(1,0)};
\addplot[gray,dashed,samples=2] coordinates {(0,0)(0,2)};
\addplot[gray,dashed,samples=2] coordinates {(1,0)(1,1)};
\node[below left] at (axis cs:0,-0.05) {$x_0=0$};
\node[below] at (axis cs:1,-0.05) {$x_1=1$};
\end{axis}
\end{tikzpicture}
\caption{A two-cycle for $f(x)=x^3-2x+2$ started at $x_0=0$. The red tangent at
$0$ crosses at $1$; the green tangent at $1$ crosses back at $0$. Newton's method
ping-pongs $0\to1\to0\to1\to\cdots$ and never finds the actual root near $-1.77$.}
\label{ch19:fig-cycle}
\end{figure}

\subsection{Failure 4: a double root slows everything to a crawl}

The final ``failure'' is gentler -- the method still converges, but it loses its
trademark speed. This happens at a \emph{multiple root}, where $f(r) = 0$ and also
$f'(r) = 0$, so the curve does not cross the axis cleanly but kisses it tangentially.
Take
\[
  f(x) = (x-1)^2, \qquad r = 1 \text{ (a double root)}.
\]
The iteration becomes
\[
  x_{n+1} = x_n - \frac{(x_n-1)^2}{2(x_n-1)} = x_n - \frac{x_n - 1}{2}
          = \frac{x_n + 1}{2}.
\]
Starting at $x_0 = 2$ this gives $1.5,\ 1.25,\ 1.125,\ 1.0625,\ \ldots$ -- the error
is merely \emph{halving} each step. That is steady but slow \textbf{linear}
convergence: one extra correct bit per step, not a doubling of digits. The
near-flat crossing starves the denominator $f'(x_n)$, which is heading toward zero
along with the numerator, and quadratic convergence is lost (Figure~\ref{ch19:fig-doubleroot}).
(There are repaired versions of Newton's method that restore fast convergence at known
multiple roots, but the plain method crawls.)

\begin{figure}[H]
\centering
\begin{tikzpicture}
\begin{axis}[
  width=11cm, height=6.5cm,
  axis lines=middle,
  xlabel={$x$}, ylabel={$y$},
  xmin=0.6, xmax=2.4, ymin=-0.3, ymax=1.5,
  xtick={\empty}, ytick={\empty},
  clip=false,
]
\addplot[curveblue, very thick, samples=100, domain=0.7:2.2] {(x-1)^2};
\addplot[tangentred, thick, samples=2, domain=1.35:2.1] {1 + 2*(x-2)};
\addplot[targetgreen, thick, samples=2, domain=1.15:1.65] {0.25 + 1.0*(x-1.5)};
\addplot[only marks, mark=*, mark size=1.5pt, black] coordinates {(2,0) (2,1) (1.5,0) (1.5,0.25) (1.25,0) (1,0)};
\addplot[gray, dashed] coordinates {(2,0) (2,1)};
\addplot[gray, dashed] coordinates {(1.5,0) (1.5,0.25)};
\node[below] at (axis cs:2,-0.05) {$x_0$};
\node[below] at (axis cs:1.5,-0.05) {$x_1$};
\node[below] at (axis cs:1.25,-0.05) {$x_2$};
\node[below left] at (axis cs:1,-0.02) {$r=1$};
\node[above left] at (axis cs:2,1) {$(x_0,f(x_0))$};
\end{axis}
\end{tikzpicture}
\caption{Why a double root crawls on $f(x)=(x-1)^2$. Because the parabola kisses the
horizontal axis at $r=1$ rather than crossing it at a slant, the tangent line from any
guess $x_n$ lands at the exact midpoint between $x_n$ and the root:
$x_{n+1}-1 = \tfrac12(x_n-1)$. Instead of squaring the error, each step merely cuts
the remaining distance in half (linear convergence).}
\label{ch19:fig-doubleroot}
\end{figure}

\begin{intuitionbox}
Collect the lesson. Newton's method is a \emph{local} method: near a simple root,
with an appropriate initial guess, it converges rapidly; but a
flat derivative, a distant start, an unlucky cycle, or a multiple root can stall,
mislead, or diverge it. The practical approach is to start from a reasonable guess
(often found by a crude, slow, but \emph{reliable} method first), then apply
Newton's method to compute the root to full precision, and \emph{always check the answer}
by evaluating $f$.
\end{intuitionbox}

\begin{historybox}
The algorithm bears Isaac Newton's name, but its evolution is a striking example of how modern calculus cleaned up historical algebra. In 1669, Newton wrote a manuscript titled \emph{De analysi per aequationes numero terminorum infinitas} (circulated privately to colleagues and published only in 1711). In it, Newton demonstrated his method on a single cubic polynomial equation:
\[
  y^3 - 2y - 5 = 0.
\]
Newton noted that $y=2$ is a close guess (giving $8 - 4 - 5 = -1$). Setting $y = 2 + p$, he expanded $(2+p)^3 - 2(2+p) - 5 = -1 + 10p + 6p^2 + p^3 = 0$. Neglecting the nonlinear terms $6p^2$ and $p^3$ as small, he solved $-1 + 10p \approx 0 \implies p = 0.1$. He then set $p = 0.1 + q$, substituted again, and repeated the process. Newton's method was purely algebraic, required expanding polynomials anew at every single stage, and made zero reference to derivatives, slopes, or tangent lines!

In 1690, \textbf{Joseph Raphson} published \emph{Analysis aequationum universalis}. Raphson recognized that Newton's tedious chain of changing variables ($y \to p \to q$) could be replaced by a clean, unified algebraic iteration: generating each new numerical approximation directly from the original polynomial coefficients. Because Raphson published this streamlined iteration fifty years before Newton's original text appeared in print, the algorithm is properly called the \textbf{Newton--Raphson method}, and priority disputes simmered for decades over historical credit.

The crowning step belongs to \textbf{Thomas Simpson}. In 1740, in \emph{Essays on Several Curious and Useful Subjects in Speculative and Mixed Mathematicks}, Simpson broke free of polynomials. He was the first to frame the method using the derivative in the denominator---$x_{n+1} = x_n - f(x_n)/f'(x_n)$---and the first to recognize that the tangent line to \emph{any} smooth geometric curve provides the path to the axis. Simpson also generalized the method to systems of simultaneous non-linear equations. It took seventy years for Newton's polynomial scratchpad to become the geometric, calculus-powered engine we use today.

Yet even with calculus in hand, fundamental mysteries remained open. In 1879, the British mathematician \textbf{Arthur Cayley} asked a deceptively simple question: if you run Newton's method on the complex cubic equation $z^3 - 1 = 0$ (whose three roots form an equilateral triangle in the complex numbers), which initial guesses converge to which root? Cayley was baffled to find that near the boundaries between the three basins of attraction---the regions of initial guesses flowing to each root---the behavior was wildly chaotic. It was only in the 1920s that Pierre Fatou and Gaston Julia proved these boundary zones are \emph{fractals}: geometric sets of infinite, self-similar detail where an infinitesimal nudge in your starting guess changes the destination completely. Today, finding provably fast starting points for general systems of polynomial equations remains a celebrated open frontier, highlighted in Steve Smale's 17th problem for twenty-first-century mathematics.
\end{historybox}

\section{Why it is so fast: quadratic convergence}
\label{ch19:quadratic}

We have repeatedly claimed that, in the happy case, the number of correct digits
\emph{doubles} each step. Time to prove it. The proof is short, and it pinpoints
exactly which features of $f$ make the method fly.

The tool is Taylor's theorem to second order, which you can think of for now as
linear approximation \emph{with its error term written out exactly}. It says: for a function
with a continuous second derivative, expanding $f$ about the point $x_n$ and
evaluating at the true root $r$,
\[
  f(r) = f(x_n) + f'(x_n)\,(r - x_n) + \tfrac12 f''(\xi)\,(r - x_n)^2,
\]
where $\xi$ is some point between $x_n$ and $r$. The left side is $f(r) = 0$,
because $r$ is the root. So
\[
  0 = f(x_n) + f'(x_n)\,(r - x_n) + \tfrac12 f''(\xi)\,(r - x_n)^2.
\]
Now divide the whole line by $f'(x_n)$ (legal, since we are assuming a simple root,
so $f'$ is not zero nearby):
\[
  0 = \frac{f(x_n)}{f'(x_n)} + (r - x_n) + \frac{f''(\xi)}{2 f'(x_n)}\,(r - x_n)^2.
\]
Move the first two terms to the left-hand side, noting that $-(r - x_n) = x_n - r$:
\[
  x_n - \frac{f(x_n)}{f'(x_n)} - r = \frac{f''(\xi)}{2 f'(x_n)}\,(r - x_n)^2.
\]
The first two terms are precisely the definition of $x_{n+1}$:
\[
  \underbrace{x_n - \frac{f(x_n)}{f'(x_n)}}_{=\,x_{n+1}} - r
  = \frac{f''(\xi)}{2 f'(x_n)}\,(r - x_n)^2.
\]
Define the error at step $n$ as $e_n = x_n - r$. Then $(r - x_n)^2 = e_n^2$, and the
left side is exactly $x_{n+1} - r = e_{n+1}$. We have arrived at the master equation
of Newton's method:
\[
  \boxed{\,e_{n+1} = \frac{f''(\xi)}{2 f'(x_n)}\; e_n^{\,2}\,}.
\]

\begin{focusbox}
Read this equation slowly, for it contains the whole story. The new error
$e_{n+1}$ is proportional to the \emph{square} of the old error $e_n$. Once $e_n$ is
small -- say $0.01$ -- its square is $0.0001$, which is a hundred times smaller. The
next error is the square of \emph{that}, around $10^{-8}$, and the next around
$10^{-16}$. Each step roughly doubles the count of accurate digits. That squaring is
quadratic convergence, and it is built into the geometry: it comes from the fact
that the tangent line matches the curve's value \emph{and} slope, so the leftover
disagreement is governed by the next term, the curvature $f''$, multiplying the
square of the gap.
\end{focusbox}

\textbf{Pause and think:} Why does squaring the error double the number of correct
decimal places? Write the error in scientific notation: an error of $10^{-3}$ means
three digits after the decimal point are verified. When you square that error,
$(10^{-3})^2 = 10^{-6}$, yielding six correct places. Square again, and you reach
$10^{-12}$, giving twelve correct places! Each step doubles the negative exponent,
multiplying your precision at exponential speed.

Figure~\ref{ch19:fig-convergence} makes the contrast between quadratic and linear
convergence visible. We plot the number of correct digits against the step number
for our $\sqrt2$ run (a simple root, quadratic) and for the double-root crawl
$f(x)=(x-1)^2$ from Section~\ref{ch19:failures} (linear). The quadratic curve
bends sharply upward -- each step buys roughly twice the digits of the last -- while
the linear one climbs at a fixed, plodding pace.

\begin{figure}[H]
\centering
\begin{tikzpicture}
\begin{axis}[
  width=11cm, height=6.5cm,
  axis lines=left,
  xlabel={step $n$}, ylabel={correct decimal digits},
  xmin=0, xmax=5, ymin=0, ymax=13,
  xtick={0,1,2,3,4,5}, ytick={0,2,4,6,8,10,12},
  legend pos=north west, clip=false,
]
% quadratic: sqrt2 errors 0.586,0.0858,0.00245,2.1e-6,1.6e-12 -> -log10
\addplot[tangentred, very thick, mark=*, mark size=1.6pt]
  coordinates {(0,0.23)(1,1.07)(2,2.61)(3,5.67)(4,11.8)};
\addlegendentry{quadratic (simple root)}
% linear: double root (x-1)^2 from x0=2 errors 1,0.5,0.25,...; digits = -log10
\addplot[curveblue, very thick, mark=square*, mark size=1.4pt]
  coordinates {(0,0)(1,0.30)(2,0.60)(3,0.90)(4,1.20)};
\addlegendentry{linear (double root)}
\end{axis}
\end{tikzpicture}
\caption{Correct digits versus step number. For a simple root (red) Newton's method
gains digits at an accelerating rate -- the signature of quadratic convergence. For a
double root (blue) it gains a fixed amount per step, the slow march of linear
convergence. Same algorithm, two radically different speeds, decided entirely by
whether $f'(r)=0$.}
\label{ch19:fig-convergence}
\end{figure}

The master equation also explains every failure mode at a single glance:

\begin{itemize}
\item \textbf{The constant out front.} The error is multiplied by $C \approx
  \abs{f''/2f'}$. If the curve has a clean, steep crossing -- $f'$ large, $f''$
  modest -- then $C$ is small and convergence is extremely fast. If $f'$ is near zero
  (the flat-tangent failure), $C$ blows up and a step can hurl you away. The
  denominator $f'$ is the method's lifeline.
\item \textbf{Why ``close enough'' matters.} The relation $e_{n+1} \approx C e_n^2$
  shrinks the error only if $C\,\abs{e_n} < 1$, i.e.\ only once $\abs{e_n} < 1/C$.
  Inside that radius each step shrinks the error and the squaring takes over; outside
  it, squaring a too-large number makes things \emph{worse}. That threshold is the
  precise meaning of ``start close enough.''
\item \textbf{Why double roots crawl.} At a double root $f'(r) = 0$, so the
  derivation above breaks -- you cannot divide by $f'(x_n)$ cleanly, because it too is
  heading to zero. A more careful analysis (which we leave to the exercises) shows
  the error then merely \emph{halves} each step: linear, not quadratic. The whole
  speed advantage came from $f'(r) \ne 0$.
\end{itemize}

\begin{aestheticbox}
The very same feature that made the
tangent line a \emph{good} approximation -- that it matches the curve to first order,
leaving an error of size $(\Delta x)^2$ -- is what makes Newton's method
\emph{fast}. The quadratic error of linear approximation
turns into the quadratic \emph{convergence} that
doubles the number of accurate digits per step.
\end{aestheticbox}

\section{Newton's method as fixed-point iteration}
\label{ch19:fixed-point}

Newton's method belongs to a larger family of root-finding algorithms, the
\textbf{fixed-point iterations}, and it is a particularly clever member of that
family. Placing it there explains \emph{why} Newton converges so much faster than
its cousins, and it ties the method to a general theory worth knowing.

A fixed-point iteration is any rule of the shape
\[
  x_{n+1} = g(x_n),
\]
for some function $g$. You feed in a number, the function spits out the next number,
and you keep feeding the output back in. If the sequence settles down to a limit
$x^\star$, then taking limits on both sides forces $x^\star = g(x^\star)$: the limit
is a point that $g$ leaves \emph{fixed}. That is what a fixed point is -- an input
$g$ returns unchanged.

Newton's method is exactly such an iteration with the particular choice
\[
  g(x) = x - \frac{f(x)}{f'(x)}.
\]
And notice: at a root $r$ of $f$ we have $f(r) = 0$, so $g(r) = r - 0 = r$. The
roots of $f$ are precisely the fixed points of Newton's $g$. Hunting a root of $f$
has been quietly converted into hunting a fixed point of $g$.

\begin{connectionbox}
Fixed points are everywhere in mathematics, and the picture for finding them is the
\emph{cobweb diagram}: plot $y = g(x)$ and the line $y = x$ on the same axes; a
fixed point is where the curve meets the diagonal. The iteration is a staircase --
go vertically to the curve, horizontally to the diagonal, vertically to the curve --
and whether the staircase spirals \emph{into} the fixed point or \emph{away} from it
is governed by one number: the slope $g'(x^\star)$ at the fixed point. If
$\abs{g'(x^\star)} < 1$ the iteration is a contraction near the fixed point and
homes in; if $\abs{g'(x^\star)} > 1$ it is repelled and flies off. This is the
contraction-mapping idea, and it reaches far beyond calculus -- into differential
equations, economics (market equilibria are fixed points), and the convergence
proofs that underpin much of numerical analysis.
\end{connectionbox}

Here is what sets Newton's method apart from an ordinary fixed-point
iteration. For a generic $g$, convergence is \emph{linear}: the error shrinks by roughly
the constant factor $\abs{g'(x^\star)}$ each step -- a fixed fraction, like steadily
losing half the error every time. That is fine but slow. Newton's particular $g$ is
engineered so that its derivative is not merely less than $1$ at the root -- it is
\emph{exactly zero} there. Here is the verification. Differentiate
$g(x) = x - f(x)/f'(x)$ using the quotient rule on the fraction:
\[
  g'(x) = 1 - \frac{f'(x)\,f'(x) - f(x)\,f''(x)}{\bigl(f'(x)\bigr)^2}
        = 1 - \frac{\bigl(f'(x)\bigr)^2 - f(x)\,f''(x)}{\bigl(f'(x)\bigr)^2}
        = \frac{f(x)\,f''(x)}{\bigl(f'(x)\bigr)^2}.
\]
Evaluate at the root $r$, where $f(r) = 0$:
\[
  g'(r) = \frac{f(r)\,f''(r)}{\bigl(f'(r)\bigr)^2} = \frac{0 \cdot f''(r)}{\bigl(f'(r)\bigr)^2} = 0.
\]

\begin{whybox}
Why does $g'(r) = 0$ explain the speed? In a fixed-point iteration the error
multiplies by about $g'(x^\star)$ each step -- so if that slope were, say, $0.5$, you
would lose half the error per step (linear). Newton arranges $g'(r) = 0$, which kills
the linear term entirely. With the first-order term gone, the leading behavior of
the error is governed by the \emph{next} term in the Taylor expansion of $g$, the
one carrying $g''(r)$ and a factor of $e_n^2$. So $e_{n+1} \approx \tfrac12
g''(r)\,e_n^2$ -- error proportional to the square of the error, which is exactly the
quadratic convergence we derived directly. Newton's method succeeds because it is the
unique formulation of a fixed-point map that zeroes out the derivative at the root, and
zeroing the derivative is precisely what upgrades linear convergence to quadratic.
\end{whybox}

\begin{figure}[H]
\centering
\begin{tikzpicture}
\begin{axis}[
  width=9.5cm, height=7cm,
  axis lines=middle,
  xlabel={$x_n$}, ylabel={$x_{n+1}$},
  xlabel style={anchor=south},
  xmin=0.9, xmax=2.1, ymin=0.9, ymax=2.1,
  xtick={1,1.5,2}, ytick={1,1.5,2},
  legend style={at={(0.95,0.20)}, anchor=east, font=\small, row sep=3pt},
  clip=false,
]
\addplot[curveblue, very thick, samples=100, domain=0.95:2.05] {x/2 + 1/x};
\addlegendentry{$y=g(x)$}
\addplot[gray, thick, samples=2, domain=0.9:2.1] {x};
\addlegendentry{$y=x$}
\addplot[tangentred, thick] coordinates {
  (2,2) (2,1.5) (1.5,1.5) (1.5,1.41667) (1.41667,1.41667) (1.41667,1.41421) (1.41421,1.41421)
};
\addplot[only marks, mark=*, mark size=1.4pt, black] coordinates {(1.41421,1.41421)};
\draw[->, thick, gray] (axis cs:1.2,1.75) -- (axis cs:1.39,1.44);
\node[above, font=\small] at (axis cs:1.2,1.75) {fixed point $r=\sqrt2$};
\end{axis}
\end{tikzpicture}
\caption{Cobweb diagram for Newton's iteration on $\sqrt2$, where $g(x)=\tfrac{x}{2}
+\tfrac1x$. The staircase (red) descends from $x_0=2$ to the fixed point where
$y=g(x)$ meets the diagonal $y=x$. Because the curve $g$ is \emph{flat} (slope $0$)
right at the fixed point, the staircase collapses onto it with astonishing speed.}
\label{ch19:fig-cobweb}
\end{figure}

\section{Newton's method in real computation}
\label{ch19:computing}

Seeing the method as code is instructive, because the gap between the formula and a
robust program is exactly where the failure modes of
Section~\ref{ch19:failures} demand respect. Here is the method as a small program,
written in plain pseudocode that any language would render almost verbatim.

\begin{lstlisting}[language=Python]
def newton(f, fprime, x0, tol=1e-12, max_iter=50):
    x = x0
    for n in range(max_iter):
        fx = f(x)
        if abs(fx) < tol:          # close enough: f(x) is essentially zero
            return x
        dfx = fprime(x)
        if dfx == 0:               # flat tangent: bail out, do not divide by zero
            raise ValueError("zero derivative; Newton stalled")
        x = x - fx / dfx           # the one line that is the whole method
    raise ValueError("did not converge in max_iter steps")
\end{lstlisting}

Three things in this little program encode the wisdom of the chapter. The
\texttt{tol} test checks the answer by evaluating $f$ -- never trust the iterate
blindly. The zero-derivative guard refuses to divide by a vanishing slope. And the
\texttt{max\_iter} cap stops the loop from running forever if the start was bad and
the iterates are wandering off or cycling. The mathematics is one line; the
engineering is the safety rail around it.

\begin{applicationbox}
Newton's method drives a great deal of modern computation, far beyond solving
single equations.
\begin{itemize}
\item \textbf{Optimization and machine learning.} To minimize a function $F$, you
  look for where its derivative is zero -- a root of $F'$. Applying Newton's method
  to $F'$ gives the update $x_{n+1} = x_n - F'(x_n)/F''(x_n)$, ``Newton's method for
  optimization,'' which uses curvature ($F''$) to take far smarter steps than plain
  gradient descent. In many variables $F''$ becomes the Hessian matrix, and training
  large models leans on quasi-Newton methods (BFGS, L-BFGS) that approximate it.
\item \textbf{Systems of equations.} Replace $f$ by a vector of functions and $f'$ by
  the Jacobian matrix; the update becomes ``solve a linear system at each step.''
  This multivariable Newton is how engineers solve the giant nonlinear systems behind
  circuit simulation, fluid dynamics, and structural analysis.
\item \textbf{Inside the hardware.} As mentioned, reciprocals, square roots, and
  divisions are computed by Newton iterations chosen so that each step needs only
  multiplication -- the operation silicon does fastest.
\end{itemize}
The reason it is chosen again and again is the quadratic convergence we proved: when
each step doubles your digits, you reach full machine precision (about $16$ decimal
digits) in a mere handful of iterations.
\end{applicationbox}

\section{A fully worked, non-algebraic example}
\label{ch19:full-example}

To close, let us solve a genuine transcendental equation by hand:
\[
  \cos x = x.
\]
This asks for the place where the cosine curve crosses the diagonal line $y = x$ --
a number called the \emph{Dottie number}, with no closed form whatsoever. Rewrite it
as a root problem:
\[
  f(x) = \cos x - x = 0, \qquad f'(x) = -\sin x - 1.
\]
A quick sketch (or noting $f(0) = 1 > 0$ and $f(1) = \cos 1 - 1 \approx -0.46 < 0$)
shows the root lies between $0$ and $1$. Start at $x_0 = 1$ (radians throughout).

\begin{labbox}
The iteration is $x_{n+1} = x_n - \dfrac{\cos x_n - x_n}{-\sin x_n - 1}$. Step
through it:
\begin{align*}
  x_0 &= 1 \\
  x_1 &= 1 - \frac{\cos 1 - 1}{-\sin 1 - 1}
       = 1 - \frac{0.540302 - 1}{-0.841471 - 1}
       = 1 - \frac{-0.459698}{-1.841471}
       = 0.750364 \\
  x_2 &= 0.750364 - \frac{\cos(0.750364) - 0.750364}{-\sin(0.750364) - 1}
       = 0.739113 \\
  x_3 &= 0.739085 \\
  x_4 &= 0.739085\ldots
\end{align*}
The root is $0.7390851332\ldots$, and we have it to nine places after four steps.
\end{labbox}

Notice the digits settling: $x_1$ is right to one place, $x_2$ to three, $x_3$ to
six or seven -- the familiar doubling. From a single crude guess and four lines of
arithmetic we have pinned down, to nine decimal places, a number that no formula can
express. This is what Newton's method buys you: it turns calculus -- the
tangent line, the derivative, the second-order error -- into a machine that
manufactures the solutions algebra cannot reach.

\begin{recoverybox}
If you take one thing from this chapter, take the master equation $e_{n+1} \approx
C\, e_n^2$ and the picture behind it: a tangent line, ridden down to the axis, again
and again. The formula $x_{n+1} = x_n - f(x_n)/f'(x_n)$ is just that picture written
in symbols, and the squaring of the error is just the price the curve charges for
not being a straight line -- a price that, near a simple root, is wonderfully small.
\end{recoverybox}

\subsection{Convergence, chaos, and algorithmic roots}
\label{ch19:what-unlocked}

\begin{connectionbox}
Newton's method is our first encounter with a \textbf{nonlinear dynamical system}:
a rule where the output of one step is fed back in as the input to the next,
producing complex evolutionary behavior from a remarkably simple formula.

When conditions are ideal, that feedback loop yields extraordinary power. Within
the \emph{basin of attraction} of a simple root---the set of all starting guesses
that successfully home in on that particular root---the error squares at every
iteration ($e_{n+1} \approx C e_n^2$). The number of verified decimal digits
doubles at each evaluation, solving \emph{transcendental equations} (equations
involving sines, cosines, or exponentials that have no algebraic formulas in
radicals) to machine precision in four or five quick steps. From celestial
mechanics predicting planetary orbits to GPS triangulation and circuit simulation,
this quadratic engine forms the backbone of quantitative engineering.

Yet when pushed into its danger zones, that same feedback loop can produce chaos.
A vanishing derivative flings iterates toward infinity. Symmetric hills trap guesses
in periodic cycles. An inflection point, as in $\arctan x$ when $|x_0| > 1.3917$,
triggers runaway divergence. And in the complex plane, the boundaries between
different basins of attraction form intricate \emph{fractals}---geometric shapes
with infinite, self-repeating detail where an infinitesimal nudge to your initial
guess can radically alter which root the method ultimately finds.

These breakdowns demonstrate why raw calculus must be tempered by numerical
discipline. Modern production libraries rarely run Newton's method in isolation.
Instead, they deploy hybrid algorithms such as \emph{Brent's method}: pairing the
guaranteed, foolproof interval-halving of the Intermediate Value Theorem with the
high-speed quadratic acceleration of Newton's method. In high-dimensional
optimization, this geometric logic blossoms into \emph{quasi-Newton algorithms}
(such as BFGS), which train modern machine learning architectures by approximating
multivariable curvature without the crushing cost of computing exact second
derivatives.

\textbf{Looking ahead.} For ten chapters we have sliced curves apart into
instantaneous slopes. In Chapter~20, we turn calculus around and begin adding things
up: the definite integral reconstructs total accumulation from a changing rate.
\end{connectionbox}

\section{Exercises}
\label{ch19:exercises}

\begin{exercise}[Mental check: resting points]
What does Newton's formula produce if your starting guess $x_0$ is already an exact
root ($f(x_0)=0$)? What happens geometrically if the tangent line at your guess is
horizontal ($f'(x_0)=0$)? State in one sentence what each situation tells you about
the iteration.
\end{exercise}

\begin{exercise}[Mental check: overshooting the domain]
Consider $f(x) = \sqrt{x}$, defined for $x \ge 0$, whose only root is $x = 0$.
(a) Compute $f'(x)$ and show that Newton's update simplifies to $x_{n+1} = -x_n$.
(b) If you start at $x_0 = 4$, where does $x_1$ land? Why does the method immediately
fail on the very next step? Sketch the tangent line at $x_0=4$ to see why it
overshoots past the boundary of the domain.
\end{exercise}

\begin{exercise}
Use Newton's method to estimate $\sqrt5$ by finding the positive root of
$f(x) = x^2 - 5$. Start at $x_0 = 2$ and carry out three iterations by hand, writing
each $x_n$ to six decimal places. Compare $x_3$ to the true value $2.2360679\ldots$
and count how many digits are correct at each step.
\end{exercise}

\begin{exercise}
Derive, from the general iteration, the formula for computing the cube root of a
positive number $a$ by applying Newton's method to $f(x) = x^3 - a$. Simplify your
update into the form $x_{n+1} = \tfrac13\bigl(2x_n + a/x_n^2\bigr)$, and use it to
estimate $\sqrt[3]{2}$ starting from $x_0 = 1$ (two iterations).
\end{exercise}

\begin{exercise}
Show that applying Newton's method to $f(x) = \dfrac1x - a$ (for a positive constant
$a$) yields the update $x_{n+1} = x_n(2 - a x_n)$, which computes $1/a$ using only
multiplication and subtraction -- no division. Test it with $a = 4$ and $x_0 = 0.2$,
performing two steps, and check that you approach $0.25$.
\end{exercise}

\begin{exercise}
The equation $x^3 - 2x + 2 = 0$ produces the two-cycle $0 \to 1 \to 0 \to \cdots$ from
the start $x_0 = 0$, as shown in the chapter. (a) Verify the cycle by direct
computation. (b) The equation has a real root near $x = -1.77$. Starting instead from
$x_0 = -2$, perform three Newton iterations and confirm the iterates head toward that
root. What does this contrast teach about the choice of starting point?
\end{exercise}

\begin{exercise}
Find the unique real root of $x^3 + x - 1 = 0$ to six decimal places using Newton's
method, starting from $x_0 = 1$. Show your work for at least three iterations and
verify your answer by plugging it back into the function.
\end{exercise}

\begin{exercise}[Conceptual]
Explain in your own words, using a sketch, why Newton's method can fail when the
starting guess lands near a point where $f'(x) = 0$. Then describe geometrically what
happens to the next iterate, and why this is a genuine breakdown of the method rather
than a minor inaccuracy.
\end{exercise}

\begin{exercise}[Conceptual]
The function $g(x) = x - f(x)/f'(x)$ is the fixed-point map behind Newton's method. A
general fixed-point iteration $x_{n+1} = g(x_n)$ converges \emph{linearly} when
$0 < \abs{g'(x^\star)} < 1$. Using the formula $g'(x) = f(x)f''(x)/(f'(x))^2$ derived
in the chapter, explain why Newton's method converges \emph{faster} than a typical
fixed-point iteration, and what specifically goes right at a simple root.
\end{exercise}

\begin{exercise}[Multiple root]
For the double root of $f(x) = (x - 3)^2$, show that Newton's iteration simplifies to
$x_{n+1} = \tfrac12(x_n + 3)$. Starting from $x_0 = 5$, compute the errors
$\abs{x_n - 3}$ for $n = 0,1,2,3,4$ and confirm that they merely \emph{halve} each
step (linear convergence) rather than the digits doubling. Explain, referring to the
master equation $e_{n+1} = \tfrac{f''(\xi)}{2f'(x_n)} e_n^2$, why the quadratic
behavior is lost at a double root.
\end{exercise}

\begin{exercise}[Proof]
Let $f$ have a continuous second derivative and a simple root $r$ (so $f'(r) \ne 0$).
Starting from the second-order Taylor expansion of $f$ about $x_n$ evaluated at $r$,
derive the master equation
\[
  e_{n+1} = \frac{f''(\xi)}{2 f'(x_n)}\, e_n^2, \qquad e_n = x_n - r,
\]
for some $\xi$ between $x_n$ and $r$. Then argue that there exist a constant $M$ and a
radius $\delta > 0$ such that if $\abs{e_0} < \delta$, the errors satisfy
$\abs{e_{n+1}} \le M \abs{e_n}^2$, and conclude that the convergence is quadratic.
\end{exercise}

\begin{exercise}[Application]
To minimize a smooth function $F(x)$, one seeks a root of $F'(x) = 0$. Write down the
Newton iteration for this root problem (it will involve $F'$ and $F''$). Apply it to
$F(x) = x^4 - 3x^2 + x$ starting from $x_0 = 1$, performing two iterations, and
interpret geometrically what the resulting point represents.
\end{exercise}
